\section{Results}
\label{sec:results}
% Every number in this section comes from tables/*.tex, figures/*.pdf or a
% numbers.tex macro.

Tables~\ref{tab:main_results} and~\ref{tab:significance} give retrieval
scores and paired tests. ``Gain'' means the mean per-query difference in
nDCG@10; headline numbers use EmbeddingGemma~2, the encoder with the better
pooled-dev score for Hybrid-fixed. Our BM25 reaches
\BMndcgSciFact{} nDCG@10 on SciFact, close to the BM25 figure reported for
BEIR~\cite{thakur2021beir}, which we take as a check that indexing and
evaluation are sound.

\inputtable{main_results}
\inputtable{significance}

\subsection{Hybrid fusion against single retrievers}
Fixed-weight fusion is significantly better than the stronger single
retriever on FiQA with both encoders, and on SciFact and HotpotQA with
BGE-small (second block of Table~\ref{tab:significance}). FiQA shows that the
two sides need not be equally strong: BM25 is far behind dense there, yet
fusion still helps. With EmbeddingGemma~2, whose dense run is much stronger
(\BestNDCGSciFact{} on SciFact), fusion significantly lowers SciFact and is
level on HotpotQA (same block). On RFC, BM25 is the best single retriever,
well ahead of either dense encoder; Hybrid-fixed without reranking differs
from it by \HybridGainOverBMRFC{} and from dense by \HybridGainOverDenseRFC. Fig.~\ref{fig:alpha_sweep} shows why: SciFact and
HotpotQA improve toward the dense end, FiQA peaks in the interior, and RFC
peaks near the sparse end, so $\alpha^{*}=\AlphaStarGemma$, chosen on the SciFact
and FiQA dev queries, sits on the wrong side of the curve for RFC. At low
dense weights the RFC curve rises above BM25 alone, so the problem there is
the transferred weight, not fusion itself.

\inputfigure{alpha_sweep}{nDCG@10 of fixed-weight fusion as a function of the
dense weight $\alpha$ on each test collection (EmbeddingGemma~2, no
reranking).}

\subsection{Does query-adaptive fusion help?}
Not reliably. With EmbeddingGemma~2 and reranking, QAF+RR minus
Hybrid-fixed+RR is \QAFgainSciFact{} on SciFact (95\% CI
\QAFgainCISciFact), \QAFgainFiQA{} on FiQA (\QAFgainCIFiQA), \QAFgainHotpot{}
on HotpotQA (\QAFgainCIHotpot) and \QAFgainRFC{} on RFC (\QAFgainCIRFC,
$p=\QAFpRFC$). Only on RFC is the difference sizeable in QAF's favour, and even there the
interval includes zero. With BGE-small, QAF+RR wins narrowly on SciFact
(first block of Table~\ref{tab:significance}) and is ahead on RFC. The
per-query spread is wide on every
collection (standard deviation of nDCG@10 for QAF+RR: \NDCGsdSciFact{} on
SciFact, \NDCGsdRFC{} on RFC), so mean differences of this size are hard to
separate from noise with our query counts.

There is room for a better weight. The plateau-centre oracle, which picks
each query's weight from its own test labels, beats $\alpha^{*}$ by
\OracleHeadroomSciFact{} on SciFact and \OracleHeadroomRFC{} on RFC. It is
a headroom reference, not an upper bound, and QAF recovers little of it.

\subsection{What QAF learned}
The ridge model explains little of the variance in the oracle targets:
training $R^2$ is \QAFtrainRsq{} on \QAFnDev{} pooled dev queries with
EmbeddingGemma~2. Among the standardised coefficients, the largest
belongs to \QAFtopFeature: a confident dense retriever pushes the weight
toward the dense side. Several features push it the other way, among them the
BM25 margin, the overlap between the two lists and, less intuitively, the
dense margin. The identifier share, which motivated the method,
has a coefficient close to zero. As a result, predictions stay close to the
intercept: on the test queries the predicted weight has mean \QAFpredMean{}
and standard deviation \QAFpredSd{}, and identifier-heavy queries get almost
the same distribution as natural-language ones. In practice QAF behaves like a
constant that has been shrunk from $\alpha^{*}=\AlphaStarGemma$ toward the
middle of the range. That shift is probably what helps on RFC, whose sweep
favours the sparse side, more than any recognition of identifier queries.
The targets explain part of the difficulty: \QAFinsensitivePct\% of training queries
have the same nDCG@10 at every grid weight and therefore carry the neutral
target 0.5.

% Coefficient table (tables/qaf_coefficients.tex) omitted for space.

\subsection{Reranking and latency}
The cross-encoder is not a uniform win. Over QAF it adds large, significant
gains on HotpotQA and RFC with both encoders (third block of
Table~\ref{tab:significance}), but it hurts strong first stages:
\RerankGainSciFact{} on SciFact with EmbeddingGemma~2, and a significant loss
on FiQA with BGE-small. A general-purpose reranker over the top 50 can undo a
good dense ranking when its notion of relevance differs from the collection's
labels. EmbeddingGemma~2's dense SciFact score may also be helped by overlap
between its training data and this public benchmark, which we cannot check.

Reranking dominates cost: with EmbeddingGemma~2 it takes \RerankLatencyMs{}~ms
of the \TotalLatencyMs{}~ms mean per-query latency of QAF+RR, while the QAF
decision takes \QAFlatencyMs{}~ms.

\subsection{Statistical caveats}
Table~\ref{tab:significance} reports 28 one-sided tests without correction
for multiple comparisons. Under a Bonferroni threshold of $0.05/28$, only the
results with $p<0.0001$ would survive: Hybrid-fixed over the best single
retriever on FiQA with EmbeddingGemma~2, reranking on HotpotQA with both
encoders and on RFC with BGE-small, and EmbeddingGemma~2 over BGE-small on
HotpotQA. None of the QAF comparisons survive. FiQA scores come from a
20{,}000-passage subset that keeps every dev and test positive, so only
relative comparisons on FiQA are meaningful.

\subsection{Ablation}
Table~\ref{tab:ablation} removes one component at a time from the full
EmbeddingGemma~2 system. Dropping the reranker raises SciFact, in line with
the reranking results above, but costs RFC most of its gain. Replacing QAF by
$\alpha^{*}$ changes little: SciFact moves slightly up and RFC slightly down.
On RFC, fixed 256-word windows score clearly below structural chunks. This
comparison changes two things at once, chunk boundaries and the size of the
retrieval unit, and its relevance labels were re-derived from answer spans,
so we read it as evidence for structure-aware chunking on RFC rather than as
a clean measurement of it. Query expansion was not evaluated.

\inputtable{ablation}

\subsection{Embedding size}
Truncating EmbeddingGemma~2 hurts dense retrieval gradually. From 768 to 256
dimensions dense nDCG@10 drops by \MatryoshkaDropTwoFiftySix{} on average over
SciFact and RFC, with RFC losing more, and the loss grows at 128
(Table~\ref{tab:matryoshka}). After fusion and reranking the curve is almost
flat, so on these two collections a 256-dimensional index costs nothing
measurable in the full pipeline.
The index file shrinks less than the dimension does, because the HNSW graph,
which does not depend on the dimension, is part of it. Dense query latency
stays in the tens of milliseconds at every size (Table~\ref{tab:matryoshka})
and is dominated by query encoding.

\inputtable{matryoshka}

\subsection{Answer quality}
Answer quality does not follow retrieval quality on HotpotQA
(Table~\ref{tab:generation}). There, Dense gives the best answers and is
significantly better in F1 than Hybrid-fixed+RR, although the reranked hybrids
retrieve better. On RFC the two hybrid configurations give the best EM, F1 and
citation accuracy, ahead of Dense and BM25. QAF+RR is not significantly
better than Hybrid-fixed+RR on either collection. In the RAGAS side analysis
on HotpotQA, Dense answers are judged slightly more faithful than QAF+RR
answers (\FaithDense{} against \FaithQAF). The judge is the same 7B model that
wrote the answers, so this comparison may suffer from self-preference and we
give it little weight. A sheet of sampled answers for manual faithfulness labelling is provided
with the code; the team has not yet completed it.

\inputtable{generation}
